% problem-000350 4 冠醚镍配合物磁性
\section*{4. 冠醚镍配合物磁性}
（图：）
图 4.1

（图：）
图 4.2

图4.3是X、Y与Cs+形成的复合物的晶体结构，在此晶体结构图中，a)为沿�轴看晶胞，b)为沿�轴看晶胞，c)为在� =0的⾯上镍配合物的⼆维排列，d)为镍配合物的⼆聚体链。

（图：）

（图：）

（图：）

d)
（图：）
图 4.3

\subsection*{4-1}
给出该晶体晶胞的组成，并写出该晶体结构基元的化学式（苯并[18]冠-6⽤X表⽰，镍配合物⽤Y表⽰）。

\subsection*{4-2}
指出镍离⼦的杂化轨道类型；图4.4是平⾯正⽅形晶体场d轨道能级分裂图，把基态镍离⼦的d电⼦填⼈该图。

图 4.4

\subsection*{4-3}
已知该晶体的晶胞参数 $t a = 1 2 7 8 . 9 9 \mathrm { p m } , \ b = 1 3 0 5 . 3 9 \mathrm { p m } , \ c = 2 7 1 7 . 0 3 \mathrm { p m } , \ \alpha = 7 8 . 3 3 9 4 ^ { \circ } , \ \beta \ = 7 7 7 . 9 9 9 9 .$ 77.0109°, $\gamma = 7 0 . 6 3 5 8 ^ { \circ } .$ 。试计算该晶体的密度 $\mathrm { ( g c m ^ { - 3 } ) }$ （相对分⼦质量：X312.4；Y451.4，三斜晶胞体积 $V = a b c \sqrt { 1 - \cos 2 \alpha - \cos 2 \beta - \cos 2 \gamma + 2 \cos \alpha \cos \beta \cos \gamma } )$

\subsection*{4-4}
写出镍配合物Y的对称元素。指出超分⼦阳离⼦中Cs+的配位数，解释其配位数较⼤的可能原因。
